%!TEX root = ../thesis.tex
% ===================================================================
% Chapter: Hardware (sits between System Design and Implementation)
% Uses the same tikzset as design_chapter.tex. If design.tex already defines
% it, delete the \tikzset block below to avoid a duplicate-definition
% warning; otherwise keep it, since this chapter draws one block diagram.
% ===================================================================
\mychapter{Hardware}{Hardware}{}
\label{chap:hardware}

\tikzset{
  blk/.style={draw, rounded corners=2pt, align=center, font=\footnotesize,
              minimum width=2.6cm, minimum height=0.95cm, inner sep=3pt},
  mdl/.style={blk, fill=black!8},
  hw/.style={blk, dashed},
  st/.style={draw, rounded corners=2pt, align=center, font=\footnotesize,
             minimum width=2.7cm, minimum height=0.9cm},
  arr/.style={-{Stealth[length=2mm]}, thick},
  lbl/.style={font=\scriptsize, fill=white, inner sep=1pt}
}

Chapter~\ref{chap:design} explained why the brain and the rover are split
across two machines. This short chapter fills in the half of that split I
left out there on purpose: what the rover itself actually is, what runs on
it out of the box, and how the brain's one line, ``the brain reaches the
Pi over rosbridge on port 9090 using roslibpy,'' turns into an actual
working connection. None of this is my own contribution, the rover is an
off-the-shelf Hiwonder kit and the ROS2 stack on it is Hiwonder's, but a
reader cannot make sense of \texttt{PiRoverController} in
Chapter~\ref{chap:implementation} without knowing what it is talking to.

\mysection{The TurboPi platform}{The TurboPi Platform}
\label{sec:turbopi_platform}

The rover is a Hiwonder TurboPi, Advanced kit, built around a Raspberry
Pi~5. The only actuators on this specific unit are a four-wheel mecanum
base, which can strafe sideways as well as drive and turn, and a
two-servo pan/tilt gimbal carrying the camera. A sonar module with its own
RGB ring sits on the front of the chassis, and a row of infrared
line-detection sensors sits underneath it. There is no arm and no
gripper on this build, and I want to be explicit about that, because the
shared ROS2 codebase Hiwonder ships still launches an
\texttt{/arm\_controller/follow\_joint\_trajectory} action server and a
\texttt{servo\_controller\_msgs/ServosPosition} topic regardless, left
over from kit variants that do have one. I confirmed with the physical
unit that neither exists here, and \texttt{PiRoverController} never
touches either topic, so I am not describing a capability the rover does
not have.

Out of the box the Advanced kit does not run a lightweight Python SDK
calling motor functions directly, it runs ROS2 Humble inside a Docker
container on the Pi, and every piece of hardware, driving, the gimbal,
the sonar RGB ring, is reached through a ROS2 topic or service rather
than a function call. A single launch file, \texttt{bringup.launch.py},
starts the nodes that matter for this project: a \texttt{mecanum}
controller that turns a velocity command into wheel speeds, a
\texttt{ros\_robot\_controller} node that is the only thing on the Pi
that actually talks to the motor and servo hardware, over a serial link
to an STM32 co-processor on the expansion board, a camera node and a
small HTTP video server, and \texttt{rosbridge\_websocket}, which is the
one node the brain ever talks to.

\begin{figure}[htbp]
  \centering
  \begin{tikzpicture}[node distance=1.1cm]
    \node[blk] (brain) at (0,0) {Brain\\(laptop)};
    \node[hw] (bridge) at (5.2,1.1) {\texttt{rosbridge\_}\\\texttt{websocket}\\:9090};
    \node[hw] (video) at (5.2,-1.1) {\texttt{web\_video\_}\\\texttt{server}\\:8080};
    \node[mdl] (mec) at (9.6,1.9) {\texttt{mecanum}};
    \node[mdl] (cam) at (9.6,0.3) {\texttt{usb\_cam}};
    \node[mdl] (rrc) at (9.6,-1.5) {\texttt{ros\_robot\_}\\\texttt{controller}};
    \node[hw] (stm) at (13.4,-1.5) {STM32\\(serial)};
    \draw[arr] (brain.east) -- node[lbl,above]{roslibpy} (bridge.west);
    \draw[arr] (brain.east) -- node[lbl,below]{HTTP} (video.west);
    \draw[arr] (bridge) -- (mec);
    \draw[arr] (video) -- (cam);
    \draw[arr] (bridge.east) -- (rrc.west);
    \draw[arr] (mec) -- (rrc);
    \draw[arr] (rrc) -- (stm);
  \end{tikzpicture}
  \caption{What the brain actually reaches on the Pi. Everything above
  \texttt{rosbridge\_websocket} and \texttt{web\_video\_server} is
  Hiwonder's stock ROS2 stack, started as one unit by
  \texttt{bringup.launch.py} inside the Pi's Docker container.}
  \label{fig:hw-bridge}
\end{figure}

\mysection{The bridge the brain actually uses}{The Bridge the Brain Actually Uses}
\label{sec:hw_bridge}

\texttt{rosbridge\_websocket} is a stock ROS2 package that exposes topics
and services as plain JSON over a WebSocket, which is the only reason
\texttt{roslibpy} works at all, it is a Python client for that JSON
protocol, not a ROS2 client in the usual sense, so the brain machine
never needs ROS2, a DDS stack, or any of Hiwonder's own packages
installed on it. I picked this over having the brain run \texttt{rclpy}
directly for a practical reason rather than a principled one: the brain
already speaks WebSocket for its own \texttt{/ws/dialogue} and
\texttt{/ws/execution} endpoints (Section~\ref{sec:protocol}), so one
more asynchronous WebSocket client sitting next to those fits the
architecture it already has, instead of asking the laptop to join the
Pi's ROS2 graph over DDS discovery across what might be a flaky Wi-Fi
link.

Driving the rover means publishing a \texttt{geometry\_msgs/Twist} on the
\texttt{cmd\_vel} topic. \texttt{linear.x} and \texttt{linear.y} are
forward and strafe speed, and \texttt{angular.z} is turn rate, but none
of the three are physical units in the way the names suggest. The
\texttt{mecanum} node sums the angular term directly into the per-wheel
inverse kinematics before scaling to a motor speed, so
\texttt{angular.z} is not radians per second, and treating it as one
would have meant guessing. I settled on working ranges by testing rather
than by calculation, roughly 3 to 8 for a turn that actually moves the
rover a sensible amount, which is also the range Hiwonder's own stock
demo code uses. The gimbal is driven the same way, a
\texttt{SetPWMServoState} message on
\texttt{/ros\_robot\_controller/pwm\_servo/set\_state} naming a servo id,
1 for tilt and 2 for pan on this unit, and a target position around a
centre of 1500. Camera frames do not go through rosbridge at all. The
brain pulls them from \texttt{web\_video\_server}'s plain HTTP snapshot
endpoint on port 8080 instead of subscribing to the raw
\texttt{sensor\_msgs/Image} topic, which means the brain needs no image
codec or ROS2 message types on its side either, a single HTTP request
returns a JPEG.

\mysection{Reusing the stock line follower}{Reusing the Stock Line Follower}
\label{sec:hw_lf}

The line-follower node \texttt{\_follow\_line} supervises
(Section~\ref{ssec:pi_rover}) is Hiwonder's own
\texttt{line\_following}, not something I wrote from scratch, and I want
to record why, because at first glance it looked unusable for this
project. Running \texttt{ros2 node info} on it shows no subscribers at
all until it is actually triggered, which looks like it might be
button-gated hardware rather than something a remote client can drive.
Reading the node's source settled that: it is dormant by design, the
camera subscription is only created inside its own \texttt{enter}
service callback, and every control point, entering, exiting, starting
or stopping the drive, and setting which line colour to follow, is a
plain ROS2 service, each one callable over rosbridge exactly like a
topic publish. Rewriting that vision pipeline myself would have meant
reimplementing a LAB colour threshold, contour extraction and a PID
steering loop for no real benefit, so \texttt{PiRoverController} drives
the stock node through its service interface instead and layers the
arrival, obstacle and timeout logic in Figure~\ref{fig:flow-followline}
on top of it.

One property of the stock node shaped that supervisor directly. It
carries its own 5-second heartbeat and exits on its own if nothing pings
it, which is sensible for a node meant to be driven from a one-shot demo
script, but it means anything that wants to keep a line-follow step
alive for longer than five seconds has to service that heartbeat itself.
\texttt{\_follow\_line} does this by polling rather than letting the
heartbeat lapse, which is also why it owns a wall-clock timeout of its
own on top of the node's: a hang on the brain's side should stop the
drive, not rely on the Pi's own node to notice first.
